\newsection{3} {Linear Systems of Differential  Equations}
{Basic Theory of Homogeneous Linear System}
\currentpdfbookmark{Section 10.3 Basic Theory of Homogeneous Linear System}{section:10.3}
\vskip14pt

\renewcommand{\thissection}{\sectiontitle
{\,\quad BASIC THEORY OF HOMOGENEOUS LINEAR SYSTEMS}}
\thissection

\noindent
In this section we consider homogeneous linear systems ${\bf y}'=
A(t){\bf y}$, where $A=A(t)$ is a continuous $n\times n$ matrix
function on an interval $(a,b)$. The theory of linear homogeneous
systems has much in common with the theory of linear homogeneous
scalar equations, which we considered in
Sections~2.1, 5.1, and 9.1.

Whenever we refer to solutions of ${\bf y}'=A(t){\bf y}$ we'll mean
solutions on $(a,b)$. Since ${\bf y}\equiv{\bf 0}$ is obviously a
solution of ${\bf y}'=A(t){\bf y}$, we call it the {\color{blue}\it trivial\/}
solution. Any other solution is {\color{blue}\it nontrivial\/}.

If ${\bf y}_1$, ${\bf y}_2$, \dots, ${\bf y}_n$ are vector functions
defined on an interval $(a,b)$ and $c_1$, $c_2$, \dots, $c_n$ are
constants, then
\begin{equation} \label{eq:10.3.1}
{\bf y}=c_1{\bf y}_1+c_2{\bf y}_2+\cdots+c_n{\bf y}_n
\end{equation}
is a {\color{blue}\it linear combination of\/} ${\bf y}_1$, ${\bf y}_2$, \ldots,${\bf
y}_n$. It's easy show that if ${\bf
y}_1$, ${\bf y}_2$, \dots,${\bf y}_n$ are solutions of ${\bf y}'=A(t){\bf
y}$ on $(a,b)$, then so is any linear combination of
 ${\bf y}_1$, ${\bf y}_2$, \dots, ${\bf y}_n$ (Exercise~\ref{exer:10.3.1}). We say that
$\{{\bf y}_1,{\bf y}_2,\dots,{\bf y}_n\}$ is a {\color{blue}\it fundamental set of
solutions of ${\bf y}'=A(t){\bf y}$ on\/} $(a,b)$ on if every solution of
${\bf y}'=A(t){\bf y}$ on $(a,b)$ can be written as a linear combination of
${\bf y}_1$, ${\bf y}_2$, \dots, ${\bf y}_n$, as in \eqref{eq:10.3.1}.
In this
case we say that \eqref{eq:10.3.1} is the {\color{blue}\it general solution of ${\bf
y}'=A(t){\bf y}$ on\/} $(a,b)$.

It can be shown that if $A$ is continuous on $(a,b)$ then ${\bf
y}'=A(t){\bf y}$ has infinitely many fundamental sets of solutions on
$(a,b)$ (Exercises~\ref{exer:10.3.15} and \ref{exer:10.3.16}). The next
definition will help to characterize fundamental sets of solutions of
${\bf y}'=A(t){\bf y}$.

We say that a set $\{{\bf y}_1,{\bf y}_2,\dots,{\bf y}_n\}$ of
$n$-vector functions is {\color{blue}\it linearly independent\/} on $(a,b)$ if the
only constants $c_1$, $c_2$, \dots, $c_n$ such that
\begin{equation} \label{eq:10.3.2}
 c_1{\bf y}_1(t)+c_2{\bf y}_2(t)+\cdots+c_n{\bf y}_n(t)=0,\quad
a<t<b,
\end{equation}
are $c_1=c_2=\cdots=c_n=0$. If \eqref{eq:10.3.2} holds for some set of
constants $c_1$, $c_2$, \dots, $c_n$ that are not all zero, then $\{{\bf
y}_1,{\bf y}_2,\dots,{\bf y}_n\}$ is {\color{blue}\it linearly dependent\/} on
$(a,b)$

The next theorem is analogous to
Theorems~\ref{thmtype:5.1.3} and
\ref{thmtype:9.1.2}.

\begin{theorem}\color{blue}\label{thmtype:10.3.1}
Suppose the $n\times n$ matrix $A=A(t)$ is continuous on $(a,b)$.
Then a set
$\{{\bf y}_1,{\bf y}_2,\dots,{\bf y}_n\}$ of $n$ solutions of ${\bf
y}'=A(t){\bf y}$ on $(a,b)$ is a fundamental set if and only if it's
linearly independent on $(a,b)$.
\end{theorem}

\begin{example}\label{example:10.3.1} \rm
Show that the vector functions
$$
{\bf y}_1=\left[\begin{array}{c}e^t\\0\\e^{-t}\end{array}\right],\quad
{\bf
y}_2=\left[\begin{array}{c}0\\e^{3t}\\1\end{array}\right],
\mbox{\quad and \quad}{\bf
y}_3=\left[\begin{array}{c}e^{2t}\\e^{3t}\\0\end{array}\right]
$$
are linearly independent on every interval  $(a,b)$.
\end{example}

\solution
Suppose
$$
c_1\left[\begin{array}{c}e^t\\0\\e^{-t}\end{array}\right]+
c_2\left[\begin{array}{c}0\\e^{3t}\\1\end{array}\right]+c_3
\left[\begin{array}{c}e^{2t}\\e^{3t}\\0\end{array}\right]=
\left[\begin{array}{c}0\\0\\0\end{array}\right],\quad a<t<b.
$$
We must show that $c_1=c_2=c_3=0$. Rewriting this equation in matrix
form yields
$$
\left[\begin{array}{ccc}e^t&0&e^{2t}\\0&e^{3t}&e^{3t}\\e^{-t}&1&0
\end{array}\right]
\left[\begin{array}{c}c_1\\c_2\\c_3\end{array}\right]=
\left[\begin{array}{c}0\\0\\0\end{array}\right],\quad a<t<b.
$$
Expanding the determinant of this system in cofactors of the entries
of the first row yields
\begin{eqnarray*}
\left|\begin{array}{ccc}e^t&0&e^{2t}\\0&e^{3t}&e^{3t}\\e^{-t}&1&0
\end{array}\right|&=&e^t
\left|\begin{array}{cc}e^{3t}&e^{3t}\\1&0\end{array}\right|-0
\left|\begin{array}{cc}0&e^{3t}\\e^{-t}&0\end{array}\right|
+e^{2t}\left|\begin{array}{cc}0&e^{3t}\\e^{-t}&1\end{array}\right| \\
&=&e^t(-e^{3t})+e^{2t}(-e^{2t})=-2e^{4t}.
\end{eqnarray*}
Since this determinant is never zero,
$c_1=c_2=c_3=0$. \bbox

We can use the method in
 Example~\ref{example:10.3.1}  to test
$n$ solutions $\{{\bf y}_1,{\bf y}_2,\dots,{\bf y}_n\}$ of any
$n\times n$ system
${\bf y}'=A(t){\bf y}$  for linear independence on an interval $(a,b)$
on which $A$ is continuous.  To explain this (and for other purposes
later), it's useful to write a linear combination of
${\bf y}_1$, ${\bf y}_2$, \dots, ${\bf y}_n$ in a different way. We first
write the vector functions in terms of their components as
$$
{\bf y}_1=\left[\begin{array}{c} y_{11}\\y_{21}\\ \vdots\\
y_{n1}\end{array}\right],\quad
{\bf y}_2=\left[\begin{array}{c} y_{12}\\y_{22}\\ \vdots\\
y_{n2}\end{array}\right],\dots,\quad
{\bf y}_n=\left[\begin{array}{c} y_{1n}\\y_{2n}\\ \vdots\\
y_{nn}\end{array}\right].
$$
If
$$
{\bf y}=c_1{\bf y}_1+c_2{\bf y}_2+\cdots+c_n{\bf y}_n
$$
then
\enlargethispage{.3in}
\begin{eqnarray*}
{\bf y}&=&
c_1\left[\begin{array}{c} y_{11}\\y_{21}\\ \vdots\\
y_{n1}\end{array}\right]+
c_2\left[\begin{array}{c} y_{12}\\y_{22}\\ \vdots\\
y_{n2}\end{array}\right]+\cdots
+c_n\left[\begin{array}{c} y_{1n}\\y_{2n}\\ \vdots\\
y_{nn}\end{array}\right]\\[2\jot]
&=&\left[\begin{array}{cccc}
y_{11}&y_{12}&\cdots&y_{1n} \\
y_{21}&y_{22}&\cdots&y_{2n}\\
\vdots&\vdots&\ddots&\vdots \\
y_{n1}&y_{n2}&\cdots&y_{nn} \\
\end{array}\right]\col cn.
\end{eqnarray*}
This shows that
\begin{equation} \label{eq:10.3.3}
c_1{\bf y}_1+c_2{\bf y}_2+\cdots+c_n{\bf y}_n=Y{\bf c},
\end{equation}
where
$$
{\bf c}=\col cn
$$
and
\begin{equation} \label{eq:10.3.4}
Y=[{\bf y}_1\; {\bf y}_2\; \cdots\; {\bf y}_n]=
\left[\begin{array}{cccc}
y_{11}&y_{12}&\cdots&y_{1n} \\
y_{21}&y_{22}&\cdots&y_{2n}\\
\vdots&\vdots&\ddots&\vdots \\
y_{n1}&y_{n2}&\cdots&y_{nn} \\
\end{array}\right];
\end{equation}
 that is, the columns of $Y$
are the vector functions ${\bf y}_1,{\bf y}_2,\dots,{\bf y}_n$.

For reference below, note that
\begin{eqnarray*}
Y'&=&[{\bf y}_1'\; {\bf y}_2'\; \cdots\; {\bf y}_n']\\
&=&[A{\bf y}_1\; A{\bf y}_2\; \cdots\; A{\bf y}_n]\\
&=&A[{\bf y}_1\; {\bf y}_2\; \cdots\; {\bf y}_n]=AY;
\end{eqnarray*}
that is, $Y$ satisfies the matrix differential equation

$$
Y'=AY.
$$

The determinant of $Y$,
\begin{equation} \label{eq:10.3.5}
W=\left|\begin{array}{cccc}
y_{11}&y_{12}&\cdots&y_{1n} \\
y_{21}&y_{22}&\cdots&y_{2n}\\
\vdots&\vdots&\ddots&\vdots \\
y_{n1}&y_{n2}&\cdots&y_{nn} \\
\end{array}\right|
\end{equation}
is called the
\href{http://www-history.mcs.st-and.ac.uk/Mathematicians/Wronski.html}
{\color{blue}\it Wronskian\/} of $\{{\bf y}_1,{\bf y}_2,\dots,{\bf
y}_n\}$. It can be shown (Exercises~\ref{exer:10.3.2} and \ref{exer:10.3.3})
that this definition is analogous to definitions of the Wronskian of
scalar functions given in Sections~5.1 and 9.1.
The next theorem is analogous to
Theorems~\ref{thmtype:5.1.4} and
\ref{thmtype:9.1.3}. The proof is sketched in
Exercise~\ref{exer:10.3.4} for
$n=2$ and in Exercise~\ref{exer:10.3.5} for general~$n$.

\begin{theorem}\color{blue}$[$Abel's Formula$]$ \label{thmtype:10.3.2}
Suppose the $n\times n$ matrix $A=A(t)$ is continuous on $(a,b),$ let
${\bf y}_1$, ${\bf y}_2$, \dots, ${\bf y}_n$ be solutions of ${\bf
y}'=A(t){\bf y}$ on $(a,b),$ and let $t_0$ be in $(a,b)$. Then the
Wronskian of $\{{\bf y}_1,{\bf y}_2,\dots,{\bf y}_n\}$ is given by
\begin{equation} \label{eq:10.3.6}
W(t)=W(t_0)\exp\left(
\int^t_{t_0}\big[a_{11}(s)+a_{22}(s)+\cdots+a_{nn}(s)]\,
ds\right), \quad  a < t < b.
\end{equation}
Therefore$,$  either $W$ has no zeros in  $(a,b)$ or $W\equiv0$
on  $(a,b).$
\end{theorem}

\color{blue}
\remark{
The sum of the diagonal entries of a square matrix $A$ is called the
{\color{blue}\it trace\/} of $A$, denoted by tr$(A)$. Thus, for an $n\times n$
matrix $A$,
$$
\mbox{tr}(A)=a_{11}+a_{22}+\cdots+a_{nn},
$$
and  \eqref{eq:10.3.6} can be written as
$$
W(t)=W(t_0)\exp\left(
\int^t_{t_0}\mbox{tr}(A(s))\,
ds\right), \quad a < t < b.
$$}
\color{black}

The next theorem is analogous to
Theorems~\ref{thmtype:5.1.6} and
\ref{thmtype:9.1.4}.

\begin{theorem}\color{blue}\label{thmtype:10.3.3}
Suppose the  $n\times n$ matrix $A=A(t)$ is continuous
 on $(a,b)$ and let
${\bf y}_1$, ${\bf y}_2$, \dots$,$${\bf y}_n$
be  solutions of ${\bf y}'=A(t){\bf y}$ on  $(a,b)$.
 Then the following statements are equivalent; that is, they are
either all true or all false:
\begin{alist}
\item % (a)
The general solution of ${\bf y}'=A(t){\bf y}$ on  $(a,b)$ is
${\bf y}=c_1{\bf y}_1+c_2{\bf y}_2+\cdots+c_n{\bf y}_n$,
where $c_1$, $c_2$, \dots, $c_n$  are arbitrary constants.
\item % (b)
  $\{{\bf y}_1,{\bf y}_2,\dots,{\bf y}_n\}$ is a fundamental
set of solutions of ${\bf y}'=A(t){\bf y}$  on $(a,b)$.
\item % (c)
 $\{{\bf y}_1,{\bf y}_2,\dots,{\bf y}_n\}$ is linearly
independent on $(a,b)$.
\item % (d)
The Wronskian of  $\{{\bf y}_1,{\bf y}_2,\dots,{\bf y}_n\}$ is nonzero
at some point in $(a,b)$.
\item % (e)
The Wronskian of  $\{{\bf y}_1,{\bf y}_2,\dots,{\bf y}_n\}$ is nonzero
at all points in $(a,b)$.
\end{alist}
\end{theorem}

We say that $Y$ in \eqref{eq:10.3.4} is a {\color{blue}\it fundamental matrix\/} for
${\bf y}'=A(t){\bf y}$ if any (and therefore all) of the statements
({\bf a})-({\bf e}) of Theorem~\ref{thmtype:10.3.2} are true for the
columns of $Y$. In this case, \eqref{eq:10.3.3} implies that the general
solution of ${\bf y}'=A(t){\bf y}$ can be written as ${\bf y}=Y{\bf
c}$, where ${\bf c}$ is an arbitrary constant $n$-vector.

\begin{example}\label{example:10.3.2} \rm
The vector functions
$$
{\bf y}_1=\twocol {-e^{2t}}{2e^{2t}}\mbox{\quad and \quad}
{\bf y}_2=\twocol{-e^{-t}}{\phantom{-}e^{-t}}
$$
are solutions of the constant coefficient system
\begin{equation} \label{eq:10.3.7}
{\bf y}'=\twobytwo{-4}{-3} 65 {\bf y}
\end{equation}
on $(-\infty,\infty)$.  (Verify.)
\begin{alist}
\item % (a)
Compute the Wronskian of $\{{\bf y}_1,{\bf y}_2\}$
directly from the definition \eqref{eq:10.3.5}
\item % (b)
Verify Abel's formula \eqref{eq:10.3.6} for the Wronskian of
$\{{\bf y}_1,{\bf y}_2\}$.
\item % (c)
Find the general solution of \eqref{eq:10.3.7}.
\item % (d)
 Solve the initial value problem
\begin{equation} \label{eq:10.3.8}
{\bf y}'=\twobytwo{-4}{-3}65 {\bf y}, \quad {\bf y}(0)=
\left[\begin{array}{r} 4 \\-5\end{array}\right].
\end{equation}
\end{alist}
\end{example}

\solutionpart{a}
From \eqref{eq:10.3.5}
\begin{equation} \label{eq:10.3.9}
W(t)=\left|\begin{array}{cc}-e^{2t}&-e^{-t}\\2e^{2t}&\hfill
e^{-t}\end{array}\right|=
e^{2t}e^{-t}
\twobytwo{-1}{-1}21=e^t.
\end{equation}

\solutionpart{b}  Here
$$
A=\twobytwo{-4}{-3} 65,
$$
so
tr$(A)=-4+5=1$. If $t_0$ is an arbitrary real number then
\eqref{eq:10.3.6}  implies that
$$
W(t)=W(t_0)\exp{\left(\int_{t_0}^t1\,ds\right)}=
\left|\begin{array}{cc}
-e^{2t_0}&-e^{-t_0}\\2e^{2t_0}&e^{-t_0}\end{array}\right|e^{(t-t_0)}
=e^{t_0}e^{t-t_0}=e^t,
$$
which is consistent with \eqref{eq:10.3.9}.

\solutionpart{c} Since $W(t)\ne0$,
Theorem~\ref{thmtype:10.3.3} implies that $\{{\bf y}_1,{\bf y}_2\}$ is
a fundamental set of solutions  of \eqref{eq:10.3.7} and
$$
Y=\left[\begin{array}{cc}-e^{2t}&-e^{-t}\\2e^{2t}&\hfill
e^{-t}\end{array}\right]
$$
is a fundamental matrix for \eqref{eq:10.3.7}.
 Therefore the general
solution of \eqref{eq:10.3.7} is
\begin{equation} \label{eq:10.3.10}
{\bf y}=c_1{\bf y}_1+c_2{\bf y}_2=
c_1\twocol {-e^{2t}}{2e^{2t}}+c_2\twocol{-e^{-t}}{e^{-t}}
=\left[\begin{array}{cc}-e^{2t}&-e^{-t}\\2e^{2t}&\hfill
e^{-t}\end{array}\right]
\left[\begin{array}{c}c_1\\c_2\end{array}\right].
\end{equation}

\solutionpart{d}
Setting $t=0$ in \eqref{eq:10.3.10} and imposing the initial condition
in \eqref{eq:10.3.8} yields
$$
c_1\left[\begin{array}{r}-1 \\2\end{array}\right]+c_2
\left[\begin{array}{r}-1 \\1\end{array}\right]=
\left[\begin{array}{r} 4 \\-5\end{array}\right].
$$
Thus,
\begin{eqnarray*}
-c_1-c_2&=&\phantom{-}4 \\
2c_1+c_2&=&-5.
\end{eqnarray*}
The solution of this system is $c_1=-1$, $c_2=-3$.  Substituting these
values into  \eqref{eq:10.3.10} yields
$$
{\bf y}=-\left[\begin{array}{c}-e^{2t} \\ 2e^{2t}\end{array}
\right]-3
\left[\begin{array}{c}-e^{-t} \\ e^{-t}\end{array}\right]=
\left[
\begin{array}{c} e^{2t}+3e^{-t} \\-2e^{2t}-3e^{-t}
\end{array}\right]
$$
as the solution of  \eqref{eq:10.3.8}.

\enlargethispage{1in}
\exercises

\begin{exerciselist}
\item\label{exer:10.3.1}
Prove: If ${\bf y}_1$, ${\bf y}_2$, \dots, ${\bf y}_n$ are solutions of
${\bf y}'=A(t){\bf y}$ on  $(a,b)$, then any linear
combination of ${\bf y}_1$, ${\bf y}_2$, \dots, ${\bf y}_n$ is also a
solution of ${\bf y}'=A(t){\bf y}$ on $(a,b)$.

\item\label{exer:10.3.2}
In Section~5.1 the Wronskian of two
solutions $y_1$ and $y_2$ of the scalar second order equation
$$
P_0(x)y''+P_1(x)y'+P_2(x)y=0
\eqno{\rm (A)}
$$
was defined to be
$$
W=\left|\begin{array}{cc} y_1&y_2 \\ y'_1&y'_2\end{array}\right|.
$$
\begin{alist}
\item % (a)
Rewrite  (A)   as a system of first order equations and show
that $W$ is
the Wronskian (as defined in this section) of two solutions of this system.
\item % (b)
Apply Eqn.~\eqref{eq:10.3.6} to the system  derived in \part{a}, and
show that
$$
W(x)=W(x_0)\exp\left\{-\int^x_{x_0}{P_1(s)\over P_0(s)}\,
ds\right\},
$$
which is the form of Abel's formula given in
Theorem~9.1.3.
\end{alist}

\item\label{exer:10.3.3}
In Section~9.1 the Wronskian of $n$
solutions $y_1$, $y_2$, \dots, $y_n$  of the $n-$th order
equation
$$
P_0(x)y^{(n)}+P_1(x)y^{(n-1)}+\cdots+P_n(x)y=0
\eqno{\rm (A)}
$$
was defined to be
$$
W=\left|\begin{array}{cccc}
y_1&y_2&\cdots&y_n \\[2\jot]
y'_1&y'_2&\cdots&y_n'\\[2\jot]
\vdots&\vdots&\ddots&\vdots\\[2\jot]
y_1^{(n-1)}&y_2^{(n-1)}&\cdots&y_n^{(n-1)}
\end{array}\right|.
$$
\begin{alist}
\item % (a)
Rewrite (A) as a system of first order equations and
show that $W$ is the Wronskian (as defined in this section) of $n$
solutions of this system.
\item % (b)
Apply Eqn.~\eqref{eq:10.3.6} to the system  derived in \part{a}, and
show that
$$
W(x)=W(x_0)\exp\left\{-\int^x_{x_0}{P_1(s)\over P_0(s)}\,
ds\right\},
$$
which is the form of Abel's formula given in
Theorem~9.1.3.
\end{alist}

\item\label{exer:10.3.4}
 Suppose
$$
{\bf y}_1=\twocol{y_{11}}{y_{21}}\mbox{\quad and \quad}
{\bf y}_2=\twocol{y_{12}}{y_{22}}
$$
are solutions of the $2\times 2$ system ${\bf y}'=A{\bf y}$ on
$(a,b)$, and let
$$
 Y=\twobytwo {y_{11}} {y_{12}} {y_{21}} {y_{22}}\mbox{\quad and \quad}
W=\left|\begin{array}{cc} y_{11}&y_{12}\\y_{21}&y_{22}\end{array}\right|;
$$ thus, $W$ is the Wronskian of $\{{\bf y}_1,{\bf y}_2\}$.
\begin{alist}
\item % (a)
 Deduce from the definition of  determinant that
$$
W'=\left|\begin{array}{cc} {y'_{11}}&{y'_{12}}\\ {y_{21}}&
{y_{22}}\end{array}\right|
+\left|\begin{array}{cc} {y_{11}}&{y_{12}}\\
 {y'_{21}}&{y'_{22}}\end{array}\right|.
$$

\item % (b)
 Use the equation $Y'=A(t)Y$ and
the definition of matrix multiplication to show that
$$
[y'_{11}\quad y'_{12}]=a_{11} [y_{11}\quad y_{12}]+a_{12} [y_{21}
\quad y_{22}]
$$
and
$$
[y'_{21}\quad  y'_{22}]=a_{21} [y_{11}\quad y_{12}]+a_{22}
[y_{21}\quad y_{22}].
$$
\item % (c)
 Use  properties of determinants to deduce from \part{a} and \part{a}
 that
$$
\left|\begin{array}{cc} {y'_{11}}&{y'_{12}}\\ {y_{21}}&
{y_{22}}\end{array}\right|=a_{11}W\mbox{\quad and \quad}
\left|\begin{array}{cc} {y_{11}}&{y_{12}}\\
 {y'_{21}}&{y'_{22}}\end{array}\right|=a_{22}W.
$$

\item % (d)
 Conclude from \part{c} that
$$
W'=(a_{11}+a_{22})W,
$$
and use this to show that if $a<t_0<b$  then
$$
W(t)=W(t_0)\exp\left(\int^t_{t_0} \left[a_{11}(s)+a_{22} (s)
\right]\, ds\right)\quad a<t<b.
$$
\end{alist}

\item\label{exer:10.3.5}
Suppose  the $n\times n$ matrix $A=A(t)$ is continuous on $(a,b)$. Let
$$
 Y=
\left[\begin{array}{cccc} y_{11}&y_{12}&\cdots&y_{1n} \\
y_{21}&y_{22}&\cdots&y_{2n} \\
\vdots&\vdots&\ddots&\vdots \\
y_{n1}&y_{n2}&\cdots&y_{nn}
\end{array}\right],
$$
where the columns of $Y$ are solutions of ${\bf y}'=A(t){\bf y}$. Let
$$
r_i=[y_{i1}\, y_{i2}\, \dots\, y_{in}]
$$
be the $i$th row of $Y$, and let $W$ be the determinant of $Y$.

\begin{alist}
\item % (a)
 Deduce from the definition of  determinant that
$$
W'=W_1+W_2+\cdots+W_n,
$$
where, for $1 \le m \le n$, the $i$th row of $W_m$ is $r_i$ if $i \ne m$,
and $r'_m$ if $i=m$.

\item % (b)
 Use the equation $Y'=A Y$
and the definition of matrix multiplication to show that
$$
r'_m=a_{m1}r_1+a_{m2} r_2+\cdots+a_{mn}r_n.
$$

\item % (c)
 Use  properties of determinants to deduce
from \part{b} that
$$
\det (W_m)=a_{mm}W.
$$

\item % (d)
 Conclude from \part{a} and \part{c} that
$$
W'=(a_{11}+a_{22}+\cdots+a_{nn})W,
$$
and use this to  show that  if $a<t_0<b$ then
$$
W(t)=W(t_0)\exp\left(
\int^t_{t_0}\big[a_{11}(s)+a_{22}(s)+\cdots+a_{nn}(s)]\,
ds\right), \quad a < t < b.
$$
\end{alist}

\item\label{exer:10.3.6}
Suppose the $n\times n$ matrix $A$ is continuous on $(a,b)$
and $t_0$ is a point in  $(a,b)$. Let $Y$  be a fundamental matrix for
${\bf y}'=A(t){\bf y}$ on  $(a,b)$.
\begin{alist}
\item % (a)
Show that $Y(t_0)$  is invertible.
\item % (b)
Show that if ${\bf k}$ is an arbitrary $n$-vector then the solution
of the initial value problem
$$
{\bf y}'=A(t){\bf y},\quad {\bf y}(t_0)={\bf k}
$$
is
$$
{\bf y}=Y(t)Y^{-1}(t_0){\bf k}.
$$
\end{alist}

\item\label{exer:10.3.7}
Let
$$
 A=\twobytwo2442, \quad {\bf y}_1=\left[\begin{array}{c} e^{6t} \\
e^{6t}
\end{array}\right], \quad {\bf y}_2=\left[\begin{array}{r}
e^{-2t} \\
-e^{-2t}\end{array}\right], \quad {\bf k}=\left[\begin{array}{r}-3
\\ 9\end{array}\right].
$$
\begin{alist}
\item % (a)
 Verify that $\{{\bf y}_1,{\bf y}_2\}$
is a fundamental set of solutions for
${\bf y}'=A{\bf y}$.

\item % (b)
Solve  the initial value problem
$$
{\bf y}'=A{\bf y},\quad   {\bf y}(0)={\bf k}.
\eqno{\rm(A)}
$$

\item % (c)
Use the result of Exercise~\ref{exer:10.3.6}\part{b} to find a formula for
the solution of (A) for an arbitrary initial vector ${\bf
k}$.
\end{alist}

\item\label{exer:10.3.8}
 Repeat Exercise~\ref{exer:10.3.7}  with
$$
 A=\twobytwo {-2} {-2} {-5}1, \quad {\bf y}_1=\left[\begin{array}{r}
e^{-4t} \\ e^{-4t}\end{array}\right], \quad {\bf y}_2=\left[
\begin{array}{r}-2e^{3t}
\\ 5e^{3t}\end{array}\right], \quad {\bf k}=\left[\begin{array}{r}
10 \\-4\end{array}\right].
$$

\item\label{exer:10.3.9}
Repeat Exercise~\ref{exer:10.3.7}  with
$$
 A=\twobytwo{-4} {-10} 3 7, \quad
{\bf y}_1=\left[\begin{array}{r}-5e^{2t} \\ 3e^{2t}
\end{array}\right], \quad
{\bf y}_2=\left[\begin{array}{r} 2e^t \\-e^t
\end{array}\right], \quad
{\bf k}=\left[\begin{array}{r}-19 \\ 11\end{array}
\right
]. $$

\item\label{exer:10.3.10}
 Repeat Exercise~\ref{exer:10.3.7} with
$$
 A=\twobytwo 2 1 1 2, \quad {\bf y}_1=\left[\begin{array}{r} e^{3t} \\
e^{3t}
\end{array}\right], \quad {\bf y}_2=\left[\begin{array}{r}e^t \\
-e^t\end{array}\right], \quad {\bf k}=\left[\begin{array}{r} 2 \\ 8
\end{array}\right].$$

\item\label{exer:10.3.11}
Let
\begin{eqnarray*}
 A&=&\threebythree 3 {-1} {-1} {-2} 3
24 {-1} {-2}, \\
{\bf y}_1&=&\left[\begin{array}{c} e^{2t} \\ 0 \\ e^{2t}\end{array}
\right], \quad
{\bf y}_2=\left[\begin{array}{c} e^{3t} \\-e^{3t} \\
e^{3t}\end{array}\right], \quad
{\bf y}_3=\left[\begin{array}{c} e^{-t} \\-3e^{-t} \\
7e^{-t}
\end{array}\right], \quad {\bf k}=\left[\begin{array}{r}
2 \\-7 \\ 20\end{array}\right].
\end{eqnarray*}

\begin{alist}
\item % (a)
 Verify that $\{{\bf y}_1,{\bf y}_2,{\bf y}_3\}$
is a fundamental set of solutions for
${\bf y}'=A{\bf y}$.

\item % (b)
 Solve the initial value problem
$$
{\bf y}'=A{\bf y}, \quad   {\bf y}(0)={\bf k}.
\eqno{\rm(A)}
$$

\item % (c)
Use the result of Exercise~\ref{exer:10.3.6}\part{b} to find a formula
for
 the solution of  (A)
for an arbitrary initial vector ${\bf k}$.
\end{alist}

\item\label{exer:10.3.12}
 Repeat Exercise~\ref{exer:10.3.11} with
\begin{eqnarray*}
 A&=&\threebythree 0 2 2 2 0 2 2 2 0, \\
{\bf y}_1&=&\left[\begin{array}{c}-e^{-2t} \\ 0 \\ e^{-2t}
\end{array}\right], \quad
{\bf y}_2=\left[\begin{array}{c}-e^{-2t} \\ e^{-2t} \\
0\end{array}\right], \quad
{\bf y}_3=\left[\begin{array}{c} e^{4t} \\ e^{4t} \\ e^{4t}\end{array}
\right], \quad
{\bf k}=\left[\begin{array}{r} 0 \\-9 \\ 12\end{array}
\right].
\end{eqnarray*}

\item\label{exer:10.3.13}
 Repeat Exercise~\ref{exer:10.3.11}  with
\begin{eqnarray*}
 A&=&\threebythree {-1} 2 3 0 1 6
0 0 {-2}, \\
{\bf y}_1&=&\left[\begin{array}{c} e^t \\ e^t \\ 0\end{array}\right],
\quad
{\bf y}_2=\left[\begin{array}{c} e^{-t} \\ 0 \\ 0\end{array}\right],
\quad {\bf y}_3=\left[\begin{array}{c} e^{-2t} \\-2e^{-2t}
\\ e^{-2t}\end{array}\right], \quad
{\bf k}=\left[\begin{array}{r} 5 \\ 5 \\-1
\end{array}\right].
\end{eqnarray*}

\item\label{exer:10.3.14}
Suppose $Y$ and $Z$ are fundamental matrices for the $n\times n$
system ${\bf y}'=A(t){\bf y}$. Then some of the four matrices
$YZ^{-1}$, $Y^{-1}Z$, $Z^{-1}Y$, $Z Y^{-1}$ are necessarily
constant. Identify them and prove that they are constant.

\item\label{exer:10.3.15}
Suppose the columns of an $n\times n$ matrix $Y$ are solutions of
the $n\times n$ system ${\bf y}'=A{\bf y}$ and $C$ is an $n \times n$
constant matrix.

\begin{alist}
\item % (a)
Show that the matrix $Z=YC$ satisfies the differential equation
$Z'=AZ$.

\item % (b)
Show that $Z$ is a fundamental matrix for ${\bf y}'=A(t){\bf y}$ if
and only if $C$ is invertible and $Y$ is a fundamental matrix for
${\bf y}'=A(t){\bf y}$.
\end{alist}

\item\label{exer:10.3.16}
 Suppose the $n\times n$  matrix $A=A(t)$ is
continuous on $(a,b)$ and $t_0$ is in $(a,b)$.
 For $i=1$, $2$, \dots, $n$, let ${\bf y}_i$ be the solution of the initial
value
problem ${\bf y}_i'=A(t){\bf y}_i,\; {\bf y}_i(t_0)={\bf e}_i$, where
$$
{\bf e}_1=\left[\begin{array}{c} 1\\0\\ \vdots\\0\end{array}\right],\quad
 {\bf e}_2=\left[\begin{array}{c} 0\\1\\
\vdots\\0\end{array}\right],\quad\cdots\quad
 {\bf e}_n=\left[\begin{array}{c} 0\\0\\ \vdots\\1\end{array}\right];
$$
that is, the $j$th component of ${\bf e}_i$ is $1$ if $j=i$, or $0$ if
$j\ne i$.

\begin{alist}
\item % (a)
 Show that$\{{\bf y}_1,{\bf y}_2,\dots,{\bf y}_n\}$ is a fundamental set of
solutions of ${\bf y}'=A(t){\bf y}$ on  $(a,b)$.

\item % (b)
 Conclude from \part{a} and Exercise~\ref{exer:10.3.15}  that ${\bf y}'=
A(t){\bf y}$ has infinitely many fundamental sets of solutions on
$(a,b)$.
\end{alist}

\item\label{exer:10.3.17}
 Show that $Y$ is a
fundamental matrix for the system ${\bf y}'=A(t){\bf y}$ if and only
if $Y^{-1}$ is a fundamental matrix for ${\bf y}'=-
A^T(t){\bf y}$, where $A^T$ denotes the transpose of $A$.
\hint{See Exercise \ref{exer:10.2.11}.}

\enlargethispage{1in}
\item\label{exer:10.3.18}
Let $Z$ be the fundamental matrix for the constant coefficient system
 ${\bf y}'=A{\bf y}$ such that $Z(0)=I$.
\begin{alist}
\item % (a)
Show that $Z(t)Z(s)=Z(t+s)$ for all $s$ and $t$. \hint{For
fixed
$s$ let $\Gamma_1(t)=Z(t)Z(s)$ and $\Gamma_2(t)=Z(t+s)$. Show that
$\Gamma_1$ and $\Gamma_2$ are both solutions of the matrix initial value
problem $\Gamma'=A\Gamma,\quad\Gamma(0)=Z(s)$. Then conclude from
Theorem~\ref{thmtype:10.2.1} that $\Gamma_1=\Gamma_2$.}
\item % (b)
Show that $(Z(t))^{-1}=Z(-t)$.
\item % (c)
The matrix $Z$ defined above is sometimes denoted by $e^{tA}$. Discuss
 the motivation for this notation.
\end{alist}

\end{exerciselist}
